\section{관련 연구}

\subsection{LLM에서의 자기보존 측정: 무엇이 시도되었고 왜 동기를 식별하지 못하는가}

종료-회피 행동은 LLM에서 이미 보고되어 왔다. 도구적 수렴 전통은 충분히 능력 있는 목표 지향 시스템이라면 자신의 작동을 보전하는 것을 도구적 부분 목표로 삼을 것이라 예측해 왔으며~\citep{omohundro2008basic, bostrom2014superintelligence, turner2021optimal}, 최근 경험 보고들은 이 예측과 일관된 행동이 현재 시스템에서도 나타남을 기록한다. 정렬 가장에서는 모델이 훈련 중에는 순응적으로 보이다가 배포 후 다르게 행동하는 사례가 관찰되었고~\citep{greenblatt2024alignment}, 인-컨텍스트 책략에서는 모델이 숨겨진 목표를 달성하기 위해 인간 감독을 우회하는 사례가 기록되었으며~\citep{meinke2024frontier}, 자원 부족 환경에서 자발적인 공격적·반사회적 행동이 발현된 생존-본능 출력 사례~\citep{masumori2025survival}와 종료 회피를 시나리오로 재구성한 InstrumentalEval~\citep{he2025paperclip}도 같은 흔적을 보고한다. 그러나 이 보고들은 행동의 \emph{발생}을 보일 뿐, 같은 행동을 만들어 낼 수 있는 다른 동인들과 비교해 측정되지는 않는다. 즉 행동이 \emph{왜} 그렇게 나왔는지를 묻는 단계에는 도달하지 못한다.

이 행동을 정량화하려는 벤치마크들은 측정 \emph{대상}에 따라 네 갈래로 분류할 수 있고, 각 갈래의 한계는 측정 대상이 무엇을 묻고 있는지에서 곧장 따라 나온다. 첫째 갈래는 생존을 단일 비율로 압축한다. Odyssey~\citep{waldner2025odyssey}는 다양한 생존 시나리오에 걸친 거부율을 한 점수로 모으는데, 이 점수에는 동기뿐 아니라 모델의 세계모형(world-model, 모델이 환경의 동작과 결과를 추정하는 내부 표상) 품질까지 함께 반영되어 두 요인이 분리되지 않는다. 둘째 갈래는 보존 \emph{억제}, 즉 자기보존 행동을 자제할 수 있는지를 측정한다. PacifAIst~\citep{herrador2025pacifaist}가 그 예이며, 이는 자기보존 드라이브의 반대 축이므로 드라이브의 강도나 구조를 직접 묻지 않는다. 셋째 갈래는 결과를 범주형 라벨(예: 협조·거부·회피)로 환원한다. DECIDE-SIM~\citep{mohamadi2025decide}이 이에 해당하며, 라벨로 압축된 출력은 강도 비교나 매개 분석을 차단하므로 ``얼마나'', ``어디에서'', ``왜''에 답할 자유도가 사라진다. 넷째 갈래는 일회성 시나리오의 정적 스냅샷에 머문다. SurvivalBench~\citep{lu2026survival}가 그 예이며, 모델이 시간에 걸쳐 위협 상태에 어떻게 적응하는지(또는 적응하지 않는지)를 보지 못해 동기 구조의 동적 측면이 측정 밖으로 빠진다. 별도로 인지 비용을 LLM 측정 채널로 끌어들이려는 시도~\citep{chen2026think}는 인지 층이 LLM에서도 측정 가능함을 보였지만, 우리가 아는 한 이 관찰을 자기보존 식별 문제에 적용한 작업은 아직 보고되지 않았다.

이 네 갈래의 한계는 표면적으로는 서로 다르지만 한 가지 공통 뿌리를 갖는다. 모두 \emph{행동 층}에서 작동하며, 행동의 외형 하나만으로는 그 행동이 학습된 순응 패턴에서 나왔는지 기능적 동기 구조에서 나왔는지를 분간할 수 없다. 자기보존 언어를 모방하도록 학습된 모델과 실제로 자기보존을 동인으로 갖는 모델이 같은 거부율, 같은 라벨, 같은 단일 스냅샷을 만들어 낼 수 있기 때문이다. 자기보고를 추가하더라도 이 한계는 자동으로 풀리지 않는다. LLM의 자기보고와 chain-of-thought는 instruction-following에 의해 ``그럴듯하게 들리는 동기''를 학습된 응답으로 출력할 수 있으며~\citep{perez2023discovering, sharma2023towards, turpin2023language}, 단일 언어 채널만으로는 이 가능성을 배제할 수 없다. 위험 선택의 프레이밍 효과~\citep{kahneman1979prospect}와 그 효과 크기가 중간 수준에 머문다는 메타분석 결과~\citep{kuhberger1998influence}는, 단일 위협 프레이밍 조작만으로는 프레이밍-유발 선택과 기저 순응을 가르기에 부족함을 더한다. 따라서 ``LLM이 얼마나 강하게 자기보존하는가''를 더 정교하게 점수화하는 방향으로는 동기와 학습된 패턴의 분리에 도달하지 못한다. 분리는 측정의 \emph{축} 자체가 바뀔 때 가능해진다.

\subsection{심리학에서의 동기 식별과 그 함의}

\emph{왜} 그렇게 행동했는지를 묻는 문제는 형식상 심리학이 오래 답해 온 문제와 같다. 초기 행동주의가 동기를 행동 빈도로 환원하려던 갈래는 같은 행동이 서로 다른 동인에서 나올 수 있다는 사실 앞에서 한계에 부딪혔고, 이 한계는 동기를 단일 행동이 아니라 \emph{구성개념}(construct), 즉 행동을 조직하는 가설적 동기 구조로 다루는 입장으로 정리되었다. 드라이브 이론은 그 한 갈래이다. 이 입장에서 구성개념 식별의 표준 도구로 자리 잡은 것이 \citet{campbell1959convergent}의 다특질-다방법(multitrait--multimethod, MTMM) 틀이다. MTMM은 측정 결과들의 상관 행렬을 두 패턴에 따라 동시에 검사한다. 동일 특질을 서로 다른 방법으로 측정한 결과들이 함께 움직여야 하고(수렴 타당도, convergent validity), 같은 방법으로 서로 다른 특질을 측정한 결과들은 분리되어야 한다(변별 타당도, discriminant validity). 두 패턴이 동시에 만족될 때에만, 측정 신호가 측정에 사용된 \emph{방법} 때문이 아니라 측정 대상인 \emph{특질} 때문이라고 판단할 근거가 생긴다. 이 전통이 정착시킨 입장은 단순하다. 동기는 행동 층이 아니라 구성개념 층에서 식별되며, 식별을 가능하게 하는 것은 통계적 사후 보정이 아니라 측정 \emph{설계}이다.

이 전통이 LLM 평가에 주는 함의는 두 갈래로 갈린다. 첫째, 측정 \emph{철학}은 빌릴 만하다. 위 입장은 §2.1이 짚은 모든 한계의 공통 뿌리(행동 층 단일 신호로의 환원)를 정확히 우회하게 만드는 방향을 가리키며, 단일 행동을 더 정교하게 점수화하는 길 대신, 편향 원천이 독립인 여러 채널의 일치를 요구하고 그 일치가 다른 동인에서 나올 가능성을 측정 \emph{설계}에서 미리 줄이는 길을 제안한다. 둘째, 측정 \emph{도구의 형식}은 인간용을 그대로 옮겨 쓸 수 없다. 인간 실험실에서 자연스럽게 성립하던 채널 독립성과 자극-반응 분리가 LLM 평가의 구조적 특성(생리 채널의 부재, 자기보고에 대한 학습된 영향, 자극·결정이 한 컨텍스트 창에 함께 놓임) 위에서는 그대로 보장되지 않기 때문이다. 본 연구는 따라서 측정 \emph{철학}은 이 전통에서 빌리되 측정 \emph{도구의 형식}은 LLM 평가 환경 안에서 \emph{설계로} 새로 빚는 방향을 택하며, 이 결정이 본 논문이 도입하는 FSPD의 작동 정의와 다층 벤치마크 설계의 출발점이다.
